\section{Two Attributes, and Only One of Them Is Forgiving}
\label{sec:ch08-two-attributes}
\index{federation!Key attribute@\code{[Key]} attribute|(}%

A subgraph needs one package and one line.

\begin{minted}{csharp}
builder.AddGraphQL()
    .AddApolloFederation()
    .AddCatalog()
\end{minted}

\code{AddApolloFederation()} does what
chapter~\ref{ch:how-a-federated-query-actually-runs} watched it do on the
sample: it adds \code{\_service} and \code{\_entities} to \code{Query}, the
\code{\_Any} scalar, an \code{\_Entity} union assembled from every type carrying
a key, and the \code{@link} that declares which version of the federation
specification the schema is written against.

What it does not do is worth saying out loud, because I expected otherwise.
Mosaic still reports the same thirteen request middleware
chapter~\ref{ch:the-life-of-a-request} printed, in the same order, with
\code{AddApolloFederation()} in place. Federation adds fields to a schema. It
does not add middleware to a pipeline. I can only say that about Mosaic:
Catalog does not carry chapter~\ref{ch:the-life-of-a-request}'s pipeline
report, so nothing in this chapter measured the new service's pipeline at all.

Then two attributes turn a class into an entity:

\begin{minted}{csharp}
[Key("id")]
public sealed record Product(
    Guid Id,
    string Sku,
    string Title,
    string? Description,
    ProductCategory Category)
{
    [ReferenceResolver]
    public static async Task<Product?> ResolveReferenceAsync(
        string id,
        IResolverContext context,
        IProductByIdDataLoader productById,
        CancellationToken cancellationToken)
        => ProductKey.TryDecode(id, context, out var productId)
            ? await productById.LoadAsync(productId, cancellationToken)
            : null;
}
\end{minted}

The key says how any other subgraph names one of these, and the reference
resolver answers when one does.
Section~\ref{sec:ch08-the-key-you-already-designed} is about the
\code{TryDecode} call in the middle, which is not boilerplate and should not
have to be there at all.

This section is about where those two attributes are written, because I put them
in the wrong place first and the schema did not complain.

\subsection*{The house style does not work here}

\index{Mosaic!authoring style}%
Every domain type in this repository is a plain record under \code{Model/},
with a partial class under \code{Types/} carrying its resolvers.
\code{Product} is the record; \code{ProductNode} is where its \code{id} field
and its node resolver live. So the obvious place for two more attributes about
\code{Product} is \code{ProductNode}, beside everything else the graph knows
about the type.

That compiles, publishes a schema and composes, and the entity it produces
cannot be resolved by anything.

\index{samples!entity-attribute-placement}%
The evidence is a sample in the companion repository,
\code{samples/entity-attribute-placement}, which declares the same trivial
entity four times, once per placement, in one schema so that the four answers
cannot be confused with each other. Here is what its \code{\_service} returns
for the type that puts both attributes on the extension class:

\begin{graphqlsdl}
type Alpha @key(fields: "id") {
  resolveByKey(id: String!): Alpha
  id: String!
  title: String!
}
\end{graphqlsdl}

\index{source generator!fields it invents}%
The key landed. The reference resolver became a public field. The source
generator collects the methods of an \code{[ObjectType<T>]} class and turns each
one into a field on the type, and it has never heard of
\code{[ReferenceResolver]}: grep the generator's source at tag \code{16.6.0},
commit \code{8fea46e}, and the only federation attribute it knows about is
\code{Shareable}, and that one belongs to the Composite Schemas stack rather
than to Apollo Federation.

So \code{resolveByKey} is now something a client can call:

\begin{minted}[fontsize=\footnotesize]{json}
{"data":{"alphas":[{"resolveByKey":{"title":"alpha one"}}]}}
\end{minted}

\index{federation!entities field@\code{\_entities} field!with no resolver}%
And \code{\_entities}, the field that actually matters, has nothing behind it:

\begin{minted}[fontsize=\footnotesize]{json}
{"errors":[{"message":"Unexpected Execution Error","path":["_entities"]}],"data":null}
\end{minted}

That is the masked error
chapter~\ref{ch:how-a-federated-query-actually-runs} taught me to read, meaning
the same thing it meant there: HotChocolate threw
\code{EntityResolver\_NoResolverFound} inside and the mask is all that survived.

\subsection*{Why it fails quietly}

\index{HotChocolate!ReferenceResolverAttribute@\code{ReferenceResolverAttribute}}%
\code{ReferenceResolverAttribute} is a descriptor attribute, and its
\code{TryConfigure} does its work only when the descriptor it is handed is an
object type descriptor or an interface type descriptor. Read at tag
\code{16.6.0}: every branch of that method is inside one of those two type
checks, and there is no \code{else}. A method inside a type extension class is
compiled into a field, so what arrives is a field descriptor, and the method
falls through and returns.

Nothing throws, nothing warns, and the failure is worse than a no-op in two
directions at once. The type still carries \code{@key}, so it is still in the
\code{\_Entity} union, so a composer will accept it and a router will route to
it. And the schema has grown a field nobody designed.

\subsection*{Where each attribute may go}

The sample's four types answer like this, against HotChocolate 16.6.0:

\begin{center}
\begin{tabular}{lllll}
\toprule
 & \code{[Key]} on & \code{[ReferenceResolver]} on & key emitted & resolves \\
\midrule
\code{Alpha}   & the extension & the extension & yes & no \\
\code{Bravo}   & the record    & the record    & yes & yes \\
\code{Charlie} & the record    & the extension & yes & no \\
\code{Delta}   & the extension & the record    & yes & yes \\
\bottomrule
\end{tabular}
\end{center}

All four are in the \code{\_Entity} union. Two of them answer a representation
and two return \enquote{Unexpected Execution Error}, and the two that fail are
exactly the two that put \code{[ReferenceResolver]} on the extension class.

\index{federation!ReferenceResolver attribute@\code{[ReferenceResolver]} attribute}%
\code{[Key]} works in either place. \code{[ReferenceResolver]} works on the
runtime type and nowhere else.

Mosaic puts both on the runtime type, and the rule for this book from here on is
that federation attributes go on the runtime type. That is an amendment to the
authoring style chapter~\ref{ch:hotchocolate-quickly} settled, not a footnote to
it, and it is recorded as one. The fields on an entity stay where the layout
rule put them; what the entity \emph{is} to other services is declared on the
class itself.

\subsection*{What a reference resolver may ask for}

\index{federation!reference resolver!parameters}%
The parameters are not ordinary resolver parameters, and the difference cost me
an afternoon. Anything the reference resolver declares is first offered to the
federation argument builder, which reads it out of the representation by name.
Measured, this is what works:

A parameter called \code{id} gets the \code{id} field of the representation.
Anything registered in the \emph{application} container arrives as a service:
\code{CatalogService}, a DataLoader, whatever the resolver needs. So does an
\code{IResolverContext}, which is how
section~\ref{sec:ch08-the-key-you-already-designed} reaches something it cannot
ask for directly.

What does not work is a service registered only with the schema. Ask for one of
those and the argument builder claims the parameter instead, looks for a
representation field by that name, finds nothing, and hands the resolver a null
that it promptly throws on. The caller sees \enquote{Unexpected Execution
Error} again.

The rule underneath is that a reference resolver's parameter list is a
description of a representation first and a dependency list second. Write it in
that order and it behaves. HotChocolate's own certification schema takes a
repository as a second parameter, so the service shape is the intended one; it
also shows a composite key spread across several parameters, with
\code{[Map]} reaching into a nested one, which Mosaic does not need and I have
not run.

\index{federation!Key attribute@\code{[Key]} attribute|)}%
